\documentclass[11pt]{article}
\usepackage[utf8]{inputenc}
\usepackage[T1]{fontenc}
\usepackage{amsmath, amssymb, amsthm}
\usepackage{geometry}
\usepackage{xcolor}
\usepackage{booktabs}
\usepackage{titlesec}
\usepackage{listings}

\geometry{a4paper, margin=1in}

\definecolor{DeepNavy}{HTML}{2E4057}
\definecolor{Teal}{HTML}{048A81}
\definecolor{WarmOrange}{HTML}{E85D04}
\definecolor{WarmGray}{HTML}{6C757D}

\lstset{
    language=Python,
    basicstyle=\ttfamily\small,
    keywordstyle=\color{DeepNavy}\bfseries,
    commentstyle=\color{WarmGray}\itshape,
    stringstyle=\color{Teal},
    backgroundcolor=\color{gray!10},
    frame=single,
    breaklines=true,
    showstringspaces=false
}

\title{\color{DeepNavy}\huge Deep Dive: Granular Instrumental Variables\\ \large Understanding the Optimal Weights $\Gamma$}
\author{Macro Lab Group — Technical Supplement}
\date{April 15, 2026}

\begin{document}

\maketitle

\section{Recap: The Simple vs. The Full Model}
In the simple model, we assume all firms respond identically to macro shocks. The GIV was simply:
\begin{equation}
z_t^{simple} = \sum (S_i - 1/N) y_{it}
\end{equation}
This works because $S_i - 1/N$ is orthogonal to a constant factor (where every firm has $\lambda_i = 1$). However, when firms have heterogeneous sensitivities $\lambda_i$, the simple weights fail. We need $\Gamma$.

\section{Intuition: $\Gamma$ as "De-biased" Firm Size}

\subsection{The Simple GIV as a Special Case}
Think of the simple GIV weights as $\Gamma_i^{simple} = S_i - \bar{S}$. We take the raw size and subtract the mean size. This "centers" the weights around zero, which kills any factor that hits everyone equally.

\subsection{The Problem with Raw Size}
If large firms are systematically more sensitive to interest rates (a common factor $\eta_t$), then the size-weighted outcome $y_{St}$ will move whenever interest rates move. If we use $y_{St}$ as an instrument, we are accidentally instrumenting with interest rate shocks, not idiosyncratic firm shocks.

\subsection{The Logic of $\Gamma$}
$\Gamma$ is the part of firm size that is \textbf{unrelated} to macro-sensitivity.
\begin{itemize}
    \item We ask: "How much of this firm's size is explained by its factor loading $\lambda_i$?"
    \item We subtract that "explained" part from the actual size.
    \item The result, $\Gamma_i$, is the "Pure Granular Weight."
\end{itemize}
If $\lambda_i$ is completely uncorrelated with $S_i$, then $\Gamma$ will look almost exactly like the simple GIV weights. If they are correlated, $\Gamma$ "tilts" the weights to protect the instrument's validity.

\section{The Math: $\Gamma$ as a Projection Residual}

Mathematically, we want to find weights $\Gamma$ that are as close as possible to the size weights $S$ (for power), but are orthogonal to the loadings $\lambda$ (for validity).

We solve the following optimization problem:
\begin{equation}
\min_{\Gamma} \sum_{i=1}^N (\Gamma_i - S_i)^2 \quad \text{subject to} \quad \sum \Gamma_i = 0, \quad \sum \Gamma_i \lambda_i = 0
\end{equation}

The solution to this problem is equivalent to the **residual** of a linear regression. If we regress the vector of sizes $S$ on a constant and the vector of loadings $\lambda$:
\begin{equation}
S_i = \alpha + \beta \lambda_i + \text{residual}_i
\end{equation}
Then the optimal weight $\Gamma_i$ is exactly that residual. 

\subsection{General Matrix Formula}
If there are multiple factors (a matrix $\Lambda$), the formula is:
\begin{equation}
\Gamma = M_{\tilde{\Lambda}} S = (I - \tilde{\Lambda} (\tilde{\Lambda}' \tilde{\Lambda})^{-1} \tilde{\Lambda}') S
\end{equation}
where $\tilde{\Lambda} = [\mathbf{1}, \Lambda]$ is the matrix of loadings plus a constant. $M_{\tilde{\Lambda}}$ is the "annihilator matrix" that projects $S$ into the space orthogonal to the factors.

\section{Python Implementation}

The following code demonstrates how to calculate $\Gamma$ using `numpy`.

\begin{lstlisting}
import numpy as np

def calculate_gamma(S, Lambda):
    """
    S: Vector of firm sizes (N x 1)
    Lambda: Matrix of factor loadings (N x K)
    """
    N = len(S)
    # Add a constant (ones) to the loadings to satisfy sum(Gamma) = 0
    ones = np.ones((N, 1))
    X = np.hstack([ones, Lambda])
    
    # Standard OLS Projection: Gamma = S - X(X'X)^-1 X'S
    # This is equivalent to the residuals of S regressed on X
    beta = np.linalg.inv(X.T @ X) @ X.T @ S
    prediction = X @ beta
    gamma = S - prediction
    
    return gamma

# Example: 5 firms, 1 factor
S = np.array([0.4, 0.3, 0.15, 0.1, 0.05])
Lambda = np.array([1.2, 0.8, 1.0, 1.1, 0.9]) # Factor loadings

gamma = calculate_gamma(S, Lambda.reshape(-1, 1))

print(f"Simple weights (S - mean): {S - 0.2}")
print(f"Optimal weights (Gamma):   {gamma}")
print(f"Sum of Gamma (should be 0): {np.sum(gamma):.4f}")
print(f"Gamma . Lambda (should be 0): {np.dot(gamma, Lambda):.4f}")
\end{lstlisting}

\section{Summary}
The transition from $z^{simple}$ to $z^{\Gamma}$ is a move from **demeaning** to **orthogonalizing**. In the simple model, we demean by subtracting $1/N$. In the full model, we orthogonalize by subtracting the projection of size onto macro-loadings.

\end{document}
